% Context from chapters-tex/fourier.tex; for mathematical reference, not compilation.
\section{Application to ODEs and PDEs}
\label{sec-fourier-pdes}

                                                
\subsection{On Continuous Domains}

We give the main calculations, leaving analytic details of existence and regularity aside.

On $\XX=\RR$ or $\TT$, one has
\eq{
	\Ff( f^{(k)} )(\om) = (\imath \om)^k \hat f(\om). 
}
Intuitively, $f^{(k)} = f \star \de^{(k)}$ where $\de^{(k)}$ is a distribution with Fourier transform $\hat \de^{(k)}(\om) = (\imath \om)^k$. 
 
Similarly on $\XX=\RR^d$ (see Section~\ref{sec-fourier-multiple-dim} for the definition of the Fourier transform in dimension $d$), one has
\eql{\label{eq-lapl-fourier}
	\Ff(\Delta f)(\om) = -\norm{\om}^2 \hat f(\om)
}
The same identity holds on $\TT^d$, with $\om$ replaced by $n \in \ZZ^d$.
 
Fourier transforms and Fourier coefficients provide a powerful tool for studying linear differential equations with constant coefficients by turning them into algebraic equations.

\begin{figure}[tbp]
\centering
\BookDrawing{tikz/fourier-heat}
\caption{Heat diffusion as a convolution.}
\label{fig-heat}
\end{figure}

As a typical example, we consider the heat equation
\eq{
	\frac{\partial f_t}{\partial t} = \Delta f_t
	\qarrq
	\foralls \om, \quad \frac{\partial \hat f_t(\om)}{\partial t} = -\norm{\om}^2 \hat f(\om).
}
Hence $\hat f_t(\om) = \hat f_0(\om) e^{-\norm{\om}^2 t}$. Applying the inverse Fourier transform and the convolution theorem gives
\eq{
	f_t = G_t \star f_0 \qwhereq
	G_t=\frac{1}{(4\pi t)^{d/2}} e^{-\frac{\norm{x}^2}{4t}}
}
On $\RR^d$, $G_t$ is a Gaussian density with standard deviation $\sqrt{2t}$ in each direction. On the torus, periodize this kernel and adjust its mass to the chosen Haar normalization.


                                                
\subsection{Finite Domain and Discretization}

On $\ZZ/N\ZZ$, corresponding to a discrete periodic domain, consider the forward first difference and the centered second difference:
\begin{align}\label{eq-disc-diff-1}
	D_1 f &\eqdef N(f_{n+1}-f_n)_n=f\star d_1\qwhereq d_1=N[-1,0,\ldots,0,1]^\top \in \RR^N, \\
	\label{eq-disc-diff-2}
	D_2 f=-D_1^\top D_1 f & \eqdef N^2 (f_{n+1}+f_{n-1}-2f_n)_n = f \star d_2 \qwhereq d_2=-d_1\star\bar d_1=N^2[-2,1,0,\ldots,0,1]^\top \in \RR^N.
\end{align}
\begin{figure}[tbp]
\centering
\BookDrawing{tikz/fourier-finite-difference-spectrum}
\caption{Comparison of the spectra of $\De$ and $D_2$.}
\label{fig-review-fourier-finite-difference-spectrum}
\end{figure}
Thanks to Proposition~\ref{prop-tfd-conv}, one can equivalently compute
\eql{\label{eq-fft-lapl}
	\Ff( D_2 f ) = \hat d_2 \odot \hat f
	\qwhereq
	(\hat d_2)_k = N^2 ( e^{\piN k}+e^{-\piN k} - 2 ) = -4N^2 \sin\pa{\frac{\pi k}{N}}^2.
}
Here $\bar d_1[n]=d_1[-n]$ denotes periodic reflection. For fixed nonzero $k$ and $N\to\infty$, one has $(\hat d_2)_k \sim -(2\pi k)^2$ which matches the scaling of~\eqref{eq-lapl-fourier}.


                                                
                                                
                                                
\section{A Bit of Group Theory}

The reference for this section is~\cite{peyre2004algebre}.

                                                
\subsection{Characters}

For $(G,+)$ a commutative group, a character is a group morphism $\chi : (G,+) \rightarrow (\CC^*,\cdot)$, i.e. it satisfies
\eq{
	\foralls (n,m) \in G, \quad \chi(n+m) = \chi(n)\chi(m).
}
The characters form the dual group $(\hat G,\odot)$ under pointwise multiplication $(\chi_1 \odot \chi_2)(n) \eqdef \chi_1(n) \chi_2(n)$. The inverse of a character $\chi$ is $\chi^{-1}(n)=\chi(-n)$.

For a finite group $G$ with $|G|=N$, we have $N \times n=0$ for every $n \in G$. Consequently, $\chi(n)^N=\chi(N n)=\chi(0)=1$: characters take values among the roots of unity,
\eql{\label{eq-char-finite-group}
	\chi(n) \in \enscond{ e^{\piN k} }{ 0 \leq k \leq N-1 }.
}
Thus $\hat G$ is a finite group (since there are finitely many maps between two finite sets) and $\chi^{-1}=\bar\chi$. In the case of a cyclic group, the dual is simple to describe.

\begin{prop}\label{prop-iso-cyclic}
	For $G=\ZZ/N\ZZ$, the dual is $\hat G=( \phi_k )_{k=0}^{N-1}$, where $\phi_k=( e^{\piN nk} )_n$. The map $k \mapsto \phi_k$ defines a noncanonical isomorphism $G \sim \hat G$.
\end{prop}
\begin{proof}
	The $\phi_k$ are indeed characters.
	
	Conversely, for any $\chi \in \hat G$, according to~\eqref{eq-char-finite-group}, $\chi(1)=e^{\piN k}$ for some $k$. Then 
	\eq{
		\chi(n